\chapter*{PHỤ LỤC B: CẤU TRÚC GÓI TIN VÀ DỮ LIỆU}
\addcontentsline{toc}{chapter}{Phụ lục B: Cấu trúc gói tin và dữ liệu}
\setcounter{section}{0}
\renewcommand{\thesection}{B\arabic{section}}

Phụ lục này trình bày cấu trúc dữ liệu thực tế được trao đổi thông qua API và cấu trúc lưu trữ trong Redis của hệ thống Secret Letter, minh chứng cho việc máy chủ hoàn toàn không nắm giữ Khóa giải mã.

% --- PHẦN B1 ---
\section*{B1. Gói tin API Tạo thông điệp (POST /api/secrets)}
\addcontentsline{toc}{section}{B1. Gói tin API Tạo thông điệp (POST /api/secrets)}

Được gửi từ Trình duyệt (Frontend) lên Máy chủ (Backend) sau khi quá trình mã hóa AES-GCM hoàn tất ở phía Client. Lưu ý rằng trường \texttt{ciphertext} chứa dữ liệu đã được mã hóa, và không hề có khóa giải mã nào được truyền đi.

\begin{minted}[language=json, caption={Payload JSON yêu cầu tạo thông điệp mã hóa}, label={lst:json-create-req}]{json}
{
  "ciphertext": "z8aX9bT...<Base64>...",
  "ttl": 86400,
  "passwordEnabled": false,
  "passwordSalt": null,
  "passwordHash": null,
  "clientMetadata": {
    "algorithm": "AES-GCM",
    "nonceLength": 12,
    "tagLength": 16
  }
}
\end{minted}

% --- PHẦN B2 ---
\section*{B2. Cấu trúc Siêu dữ liệu (Metadata) lưu trữ trên Redis}
\addcontentsline{toc}{section}{B2. Cấu trúc Siêu dữ liệu (Metadata) lưu trữ trên Redis}

Khi hệ thống nhận được yêu cầu, Backend sẽ tách biệt \texttt{ciphertext} và siêu dữ liệu (Metadata) thành 2 bản ghi riêng biệt trong Redis. Dưới đây là cấu trúc của bản ghi Metadata:

\begin{minted}[language=json, caption={Cấu trúc bản ghi Metadata trong Redis}, label={lst:json-redis-meta}]{json}
{
  "id": "sec_2aB9c8D7eF",
  "status": "active",
  "expiresAtUnix": 1729483200,
  "createdAtUnix": 1729396800,
  "passwordEnabled": false,
  "viewCount": 0,
  "clientMetadata": {
    "algorithm": "AES-GCM",
    "nonceLength": 12,
    "tagLength": 16
  }
}
\end{minted}

\textit{Ghi chú:} Bản ghi chứa \texttt{ciphertext} thực tế được lưu tại một khóa (Key) khác (ví dụ: \texttt{secret:payload:sec\_2aB9c8D7eF}) và sẽ bị xóa ngay lập tức (Atomic Delete) bằng lệnh Lua Script ở Phụ lục A khi có người đọc thành công, trong khi Metadata vẫn được giữ lại thêm một thời gian để hiển thị trạng thái "Đã bị hủy" (Burn-after-reading).